\honest{Can Humanism Match Religion's Output?}{Eliezer Yudkowsky}{2009}

Perhaps the largest voluntary institution in the world is the Catholic Church, held together by its members' gifts. It is too large to be held together by individual negotiation, as a hunter-gatherer band is. In a larger world with faster transmission, we can expect more virulent memes. It is held together by ``affective death spirals'' around its ideas, institutions and leaders, by promises of eternal happiness and damnation, and by conformity among people who meet in person. ``The Old Testament doesn't talk about Hell, but the New Testament does.'' \nb{Broadly right: the Hebrew Bible's Sheol takes the righteous and the wicked alike.} True, ``theologians don't really believe that stuff,'' but many ordinary Catholics do. \nb{No evidence is given. The linked post is about one conversation with a woman who was nominally an Orthodox Jew.}

We who call ourselves ``rationalists'' think ourselves too good for such bonds. For a simple relief project, such as food and shelter after a tidal wave in Thailand, you would do far better to ask the Pope than Richard Dawkins. So long as that is true, atheism's gains at Catholicism's expense are somewhat hollow. The Church does harm too: it opposes condoms in Africa and spends much of its money on religion. \nb{Ten days before the post, Pope Benedict XVI had said that AIDS could not be overcome by distributing condoms, ``which even aggravates the problems.''} Unclear thinking is not harmless. But avoiding harm is an empty victory if it is your only one. Perhaps the wiser but less motivated can find efficient interventions and buy good cheaply, but few of us really do that, as opposed to planning to.

So, counting only the good, does the average Catholic do more than the average atheist, just by being more active? \nb{The post does not argue the point, but American survey data from 2000 supported it for religious practice in general: weekly attenders gave more than three times as much as the rarely religious, and more of them gave and volunteered for nonreligious causes too.}

Perhaps, you may say, false beliefs about hell simply motivate more than any true belief. ``This is a fair point.'' The folk theorem that a rational agent does at least as well as an irrational one assumes it can adopt whatever policy wins. If you cannot choose unlimited mental energy, some false beliefs may motivate more than any true one, and with altruistic akrasia the God-fearing may win. But believers still sin, because we all ``run out of mental energy.'' \nb{The ego-depletion research behind that phrase later failed large replications; see ``Superstimuli and the Collapse of Western Civilization.''}

They sin and then torment themselves, as smokers reproach themselves for not quitting. And who says that caring about real people cannot move a brain as much as a heaven it cannot picture? The brain cannot visualize 3\^{}\^{}\^{}3, let alone infinity, and anything involving more than a hundred people already involves utilities too large to picture. ``Cognitive-behavioral therapy and Zen meditation are two mental disciplines experimentally shown to yield real improvements,'' so an art built on them might help. \nb{True of CBT. For meditation, a 2007 government review found mostly poor-quality studies and no firm conclusions.} If we had an evidence-based art of fighting akrasia, why must we be less motivated than a disorganized mind that fears God's wrath? That, I admit, is ``a further-future speculation.'' I offer it to show that we should not give up on rationality so fast: understanding what goes wrong, trying intelligently to fix it, and testing whether it worked is a powerful idiom.

``Really, I suspect'' the Church's power comes less from damnation than from people meeting in person. \nb{Later survey research by Robert Putnam and David Campbell fits this: religious Americans with many friends at church were likely to be more generous, and fear of God's punishment was not the reason.} Our readers are rare and scattered; if we all lived within five miles of each other, ``I bet'' we would do more, from motivation, not coordination. Perhaps videoconferencing would give some of the effect; I suspect not, but it might be worth trying.

I would like to say that fighting akrasia alone is difficult but possible, but I am not sure that is true. The largest step, I suspect, would be regular meetings, for motivation. We could also try norms that applaud caring strongly about something and expect everyone to do something useful with their life; religion does not stress getting things done. If rationalists matched even half the average Catholic's altruistic effort, better targeting could let the typical rationalist do twice as much.

The Church spends ``More than 50\%, I would venture'' on religion, so a norm of giving 5\% of income to real causes would match what a 10\% tithe gives them, and choosing causes where good is orders of magnitude cheaper would do more. \nb{Few American churchgoers tithe, and among American denominations Catholics tend to give the least.} But first, secular humanists should at least match the believers' output per head.
